\documentclass[10pt,a4paper]{amsart}
\usepackage{amsmath,amssymb,amsthm,graphicx}
% For commutative diagrams we use Xy-pic
\usepackage[all]{xy}

% We use 2cell for 2-commutative diagrams.
\xyoption{2cell}
% Theorem environments.
%
\theoremstyle{plain}
\newtheorem{theorem}[subsection]{Theorem}
\newtheorem{proposition}[subsection]{Proposition}
\newtheorem{lemma}[subsection]{Lemma}

\theoremstyle{definition}
\newtheorem{definition}[subsection]{Definition}
\newtheorem{example}[subsection]{Example}
\newtheorem{exercise}[subsection]{Exercise}
\newtheorem{situation}[subsection]{Situation}

\theoremstyle{remark}
\newtheorem{remark}[subsection]{Remark}
\newtheorem{remarks}[subsection]{Remarks}

\numberwithin{equation}{subsection}

% Macros
%
\def\lim{\mathop{\mathrm{lim}}\nolimits}
\def\colim{\mathop{\mathrm{colim}}\nolimits}
\def\Spec{\mathop{\mathrm{Spec}}}
\def\Hom{\mathop{\mathrm{Hom}}\nolimits}
\def\Ext{\mathop{\mathrm{Ext}}\nolimits}
\def\SheafHom{\mathop{\mathcal{H}\!\mathit{om}}\nolimits}
\def\SheafExt{\mathop{\mathcal{E}\!\mathit{xt}}\nolimits}
\def\Sch{\mathit{Sch}}
\def\Mor{\mathop{\mathrm{Mor}}\nolimits}
\def\Ob{\mathop{\mathrm{Ob}}\nolimits}
\def\Sh{\mathop{\mathit{Sh}}\nolimits}
\def\NL{\mathop{N\!L}\nolimits}
\def\CH{\mathop{\mathrm{CH}}\nolimits}
\def\proetale{{pro\text{-}\acute{e}tale}}
\def\etale{{\acute{e}tale}}
\def\QCoh{\mathit{QCoh}}
\def\Ker{\mathop{\mathrm{Ker}}}
\def\Im{\mathop{\mathrm{Im}}}
\def\Coker{\mathop{\mathrm{Coker}}}
\def\Coim{\mathop{\mathrm{Coim}}}

% Boxtimes
%
\DeclareMathSymbol{\boxtimes}{\mathbin}{AMSa}{"02}

%
% Macros for moduli stacks/spaces
%
\def\QCohstack{\mathcal{QC}\!\mathit{oh}}
\def\Cohstack{\mathcal{C}\!\mathit{oh}}
\def\Spacesstack{\mathcal{S}\!\mathit{paces}}
\def\Quotfunctor{\mathrm{Quot}}
\def\Hilbfunctor{\mathrm{Hilb}}
\def\Curvesstack{\mathcal{C}\!\mathit{urves}}
\def\Polarizedstack{\mathcal{P}\!\mathit{olarized}}
\def\Complexesstack{\mathcal{C}\!\mathit{omplexes}}
% \Pic is the operator that assigns to X its picard group, usage \Pic(X)
% \Picardstack_{X/B} denotes the Picard stack of X over B
% \Picardfunctor_{X/B} denotes the Picard functor of X over B
\def\Pic{\mathop{\mathrm{Pic}}\nolimits}
\def\Picardstack{\mathcal{P}\!\mathit{ic}}
\def\Picardfunctor{\mathrm{Pic}}
\def\Deformationcategory{\mathcal{D}\!\mathit{ef}}
\newenvironment{reference}{}{}
\newenvironment{slogan}{}{}
\newenvironment{history}{}{}
\usepackage{hyperref}
\usepackage{cleveref}
\crefname{originalproposition}{original proposition}{original propositions}
\Crefname{originalproposition}{Original proposition}{Original propositions}
\makeatletter
\newlabel{restricted-proposition-glue-modification}{{0GI3}{source}{}{sourcepointer-1}{}}
\newlabel{restricted-proposition-glue-modification@cref}{{[originalproposition][0][]0GI3}{[1][0][]source}{}{}{}}
\makeatother
\crefname{originalsection}{original section}{original sections}
\Crefname{originalsection}{Original section}{Original sections}
\makeatletter
\newlabel{section-noetherian}{{0GE1}{source}{}{sourcepointer-2}{}}
\newlabel{section-noetherian@cref}{{[originalsection][0][]0GE1}{[1][0][]source}{}{}{}}
\makeatother
\crefname{originalsection}{original section}{original sections}
\Crefname{originalsection}{Original section}{Original sections}
\makeatletter
\newlabel{spaces-limits-remark-limit-preserving}{{05N0}{source}{}{sourcepointer-3}{}}
\newlabel{spaces-limits-remark-limit-preserving@cref}{{[originalsection][0][]05N0}{[1][0][]source}{}{}{}}
\makeatother
\crefname{originallemma}{original lemma}{original lemmas}
\Crefname{originallemma}{Original lemma}{Original lemmas}
\makeatletter
\newlabel{lemma-canonical-extension}{{0GE3}{source}{}{sourcepointer-4}{}}
\newlabel{lemma-canonical-extension@cref}{{[originallemma][0][]0GE3}{[1][0][]source}{}{}{}}
\makeatother
\crefname{originalproposition}{original proposition}{original propositions}
\Crefname{originalproposition}{Original proposition}{Original propositions}
\makeatletter
\newlabel{more-algebra-proposition-equivalence}{{05ER}{source}{}{sourcepointer-5}{}}
\newlabel{more-algebra-proposition-equivalence@cref}{{[originalproposition][0][]05ER}{[1][0][]source}{}{}{}}
\makeatother
\crefname{originallemma}{original lemma}{original lemmas}
\Crefname{originallemma}{Original lemma}{Original lemmas}
\makeatletter
\newlabel{spaces-cohomology-lemma-proper-over-affine-cohomology-finite}{{08AS}{source}{}{sourcepointer-6}{}}
\newlabel{spaces-cohomology-lemma-proper-over-affine-cohomology-finite@cref}{{[originallemma][0][]08AS}{[1][0][]source}{}{}{}}
\makeatother
\crefname{originallemma}{original lemma}{original lemmas}
\Crefname{originallemma}{Original lemma}{Original lemmas}
\makeatletter
\newlabel{spaces-cohomology-lemma-flat-base-change-cohomology}{{073K}{source}{}{sourcepointer-7}{}}
\newlabel{spaces-cohomology-lemma-flat-base-change-cohomology@cref}{{[originallemma][0][]073K}{[1][0][]source}{}{}{}}
\makeatother
\crefname{originaltheorem}{original theorem}{original theorems}
\Crefname{originaltheorem}{Original theorem}{Original theorems}
\makeatletter
\newlabel{spaces-cohomology-theorem-formal-functions}{{08AZ}{source}{}{sourcepointer-8}{}}
\newlabel{spaces-cohomology-theorem-formal-functions@cref}{{[originaltheorem][0][]08AZ}{[1][0][]source}{}{}{}}
\makeatother
\crefname{originallemma}{original lemma}{original lemmas}
\Crefname{originallemma}{Original lemma}{Original lemmas}
\makeatletter
\newlabel{restricted-lemma-rig-surjective-alternative-permanence}{{0GH5}{source}{}{sourcepointer-9}{}}
\newlabel{restricted-lemma-rig-surjective-alternative-permanence@cref}{{[originallemma][0][]0GH5}{[1][0][]source}{}{}{}}
\makeatother
\crefname{originallemma}{original lemma}{original lemmas}
\Crefname{originallemma}{Original lemma}{Original lemmas}
\makeatletter
\newlabel{restricted-lemma-rig-etale-alternative-permanence}{{0GH4}{source}{}{sourcepointer-10}{}}
\newlabel{restricted-lemma-rig-etale-alternative-permanence@cref}{{[originallemma][0][]0GH4}{[1][0][]source}{}{}{}}
\makeatother
\crefname{originallemma}{original lemma}{original lemmas}
\Crefname{originallemma}{Original lemma}{Original lemmas}
\makeatletter
\newlabel{restricted-lemma-rig-etale-descent}{{0GH6}{source}{}{sourcepointer-11}{}}
\newlabel{restricted-lemma-rig-etale-descent@cref}{{[originallemma][0][]0GH6}{[1][0][]source}{}{}{}}
\makeatother
\crefname{originallemma}{original lemma}{original lemmas}
\Crefname{originallemma}{Original lemma}{Original lemmas}
\makeatletter
\newlabel{restricted-lemma-permanence-rig-surjective}{{0AQT}{source}{}{sourcepointer-12}{}}
\newlabel{restricted-lemma-permanence-rig-surjective@cref}{{[originallemma][0][]0AQT}{[1][0][]source}{}{}{}}
\makeatother
\crefname{originallemma}{original lemma}{original lemmas}
\Crefname{originallemma}{Original lemma}{Original lemmas}
\makeatletter
\newlabel{restricted-lemma-closed-immersion-rig-smooth-rig-surjective}{{0GD6}{source}{}{sourcepointer-13}{}}
\newlabel{restricted-lemma-closed-immersion-rig-smooth-rig-surjective@cref}{{[originallemma][0][]0GD6}{[1][0][]source}{}{}{}}
\makeatother
\crefname{originalsection}{original section}{original sections}
\Crefname{originalsection}{Original section}{Original sections}
\makeatletter
\newlabel{limits-section-limits}{{01YV}{source}{}{sourcepointer-14}{}}
\newlabel{limits-section-limits@cref}{{[originalsection][0][]01YV}{[1][0][]source}{}{}{}}
\makeatother
\crefname{originallemma}{original lemma}{original lemmas}
\Crefname{originallemma}{Original lemma}{Original lemmas}
\makeatletter
\newlabel{limits-lemma-inverse-limit-top}{{0CUF}{source}{}{sourcepointer-15}{}}
\newlabel{limits-lemma-inverse-limit-top@cref}{{[originallemma][0][]0CUF}{[1][0][]source}{}{}{}}
\makeatother
\crefname{originallemma}{original lemma}{original lemmas}
\Crefname{originallemma}{Original lemma}{Original lemmas}
\makeatletter
\newlabel{topology-lemma-describe-limits}{{0A2P}{source}{}{sourcepointer-16}{}}
\newlabel{topology-lemma-describe-limits@cref}{{[originallemma][0][]0A2P}{[1][0][]source}{}{}{}}
\makeatother
\crefname{originalproposition}{original proposition}{original propositions}
\Crefname{originalproposition}{Original proposition}{Original propositions}
\makeatletter
\newlabel{limits-proposition-characterize-locally-finite-presentation}{{01ZC}{source}{}{sourcepointer-17}{}}
\newlabel{limits-proposition-characterize-locally-finite-presentation@cref}{{[originalproposition][0][]01ZC}{[1][0][]source}{}{}{}}
\makeatother
\crefname{originallemma}{original lemma}{original lemmas}
\Crefname{originallemma}{Original lemma}{Original lemmas}
\makeatletter
\newlabel{formal-spaces-lemma-monomorphism-iso-over-red}{{0GHZ}{source}{}{sourcepointer-18}{}}
\newlabel{formal-spaces-lemma-monomorphism-iso-over-red@cref}{{[originallemma][0][]0GHZ}{[1][0][]source}{}{}{}}
\makeatother
\crefname{originallemma}{original lemma}{original lemmas}
\Crefname{originallemma}{Original lemma}{Original lemmas}
\makeatletter
\newlabel{topologies-lemma-limit-fppf-topology}{{049N}{source}{}{sourcepointer-19}{}}
\newlabel{topologies-lemma-limit-fppf-topology@cref}{{[originallemma][0][]049N}{[1][0][]source}{}{}{}}
\makeatother
\crefname{originallemma}{original lemma}{original lemmas}
\Crefname{originallemma}{Original lemma}{Original lemmas}
\makeatletter
\newlabel{limits-lemma-descend-opens}{{01Z4}{source}{}{sourcepointer-20}{}}
\newlabel{limits-lemma-descend-opens@cref}{{[originallemma][0][]01Z4}{[1][0][]source}{}{}{}}
\makeatother
\crefname{originaltheorem}{original theorem}{original theorems}
\Crefname{originaltheorem}{Original theorem}{Original theorems}
\makeatletter
\newlabel{restricted-theorem-dilatations-general}{{0ARB}{source}{}{sourcepointer-21}{}}
\newlabel{restricted-theorem-dilatations-general@cref}{{[originaltheorem][0][]0ARB}{[1][0][]source}{}{}{}}
\makeatother
\crefname{originallemma}{original lemma}{original lemmas}
\Crefname{originallemma}{Original lemma}{Original lemmas}
\makeatletter
\newlabel{restricted-lemma-output-etale}{{0GI4}{source}{}{sourcepointer-22}{}}
\newlabel{restricted-lemma-output-etale@cref}{{[originallemma][0][]0GI4}{[1][0][]source}{}{}{}}
\makeatother
\crefname{originallemma}{original lemma}{original lemmas}
\Crefname{originallemma}{Original lemma}{Original lemmas}
\makeatletter
\newlabel{restricted-lemma-output-separated}{{0ARW}{source}{}{sourcepointer-23}{}}
\newlabel{restricted-lemma-output-separated@cref}{{[originallemma][0][]0ARW}{[1][0][]source}{}{}{}}
\makeatother
\crefname{originalproposition}{original proposition}{original propositions}
\Crefname{originalproposition}{Original proposition}{Original propositions}
\makeatletter
\newlabel{spaces-morphisms-proposition-locally-quasi-finite-separated-over-scheme}{{03XX}{source}{}{sourcepointer-24}{}}
\newlabel{spaces-morphisms-proposition-locally-quasi-finite-separated-over-scheme@cref}{{[originalproposition][0][]03XX}{[1][0][]source}{}{}{}}
\makeatother
\crefname{originallemma}{original lemma}{original lemmas}
\Crefname{originallemma}{Original lemma}{Original lemmas}
\makeatletter
\newlabel{limits-lemma-descend-finite-presentation}{{01ZM}{source}{}{sourcepointer-25}{}}
\newlabel{limits-lemma-descend-finite-presentation@cref}{{[originallemma][0][]01ZM}{[1][0][]source}{}{}{}}
\makeatother
\crefname{originallemma}{original lemma}{original lemmas}
\Crefname{originallemma}{Original lemma}{Original lemmas}
\makeatletter
\newlabel{limits-lemma-eventually-proper}{{081F}{source}{}{sourcepointer-26}{}}
\newlabel{limits-lemma-eventually-proper@cref}{{[originallemma][0][]081F}{[1][0][]source}{}{}{}}
\makeatother
\crefname{originallemma}{original lemma}{original lemmas}
\Crefname{originallemma}{Original lemma}{Original lemmas}
\makeatletter
\newlabel{limits-lemma-descend-isomorphism}{{081E}{source}{}{sourcepointer-27}{}}
\newlabel{limits-lemma-descend-isomorphism@cref}{{[originallemma][0][]081E}{[1][0][]source}{}{}{}}
\makeatother
\crefname{originallemma}{original lemma}{original lemmas}
\Crefname{originallemma}{Original lemma}{Original lemmas}
\makeatletter
\newlabel{limits-lemma-descend-section}{{01Z0}{source}{}{sourcepointer-28}{}}
\newlabel{limits-lemma-descend-section@cref}{{[originallemma][0][]01Z0}{[1][0][]source}{}{}{}}
\makeatother
\crefname{originallemma}{original lemma}{original lemmas}
\Crefname{originallemma}{Original lemma}{Original lemmas}
\makeatletter
\newlabel{algebra-lemma-Noetherian-power-ideal-kills-module}{{00L6}{source}{}{sourcepointer-29}{}}
\newlabel{algebra-lemma-Noetherian-power-ideal-kills-module@cref}{{[originallemma][0][]00L6}{[1][0][]source}{}{}{}}
\makeatother
\crefname{originalsection}{original section}{original sections}
\Crefname{originalsection}{Original section}{Original sections}
\makeatletter
\newlabel{algebra-section-completion-noetherian}{{0BNH}{source}{}{sourcepointer-30}{}}
\newlabel{algebra-section-completion-noetherian@cref}{{[originalsection][0][]0BNH}{[1][0][]source}{}{}{}}
\makeatother
\crefname{originallemma}{original lemma}{original lemmas}
\Crefname{originallemma}{Original lemma}{Original lemmas}
\makeatletter
\newlabel{restricted-lemma-Noetherian-finite-type-red}{{0GC6}{source}{}{sourcepointer-31}{}}
\newlabel{restricted-lemma-Noetherian-finite-type-red@cref}{{[originallemma][0][]0GC6}{[1][0][]source}{}{}{}}
\makeatother
\crefname{originaldefinition}{original definition}{original definitions}
\Crefname{originaldefinition}{Original definition}{Original definitions}
\makeatletter
\newlabel{restricted-definition-formal-modification}{{0GDK}{source}{}{sourcepointer-32}{}}
\newlabel{restricted-definition-formal-modification@cref}{{[originaldefinition][0][]0GDK}{[1][0][]source}{}{}{}}
\makeatother
\crefname{originalsection}{original section}{original sections}
\Crefname{originalsection}{Original section}{Original sections}
\makeatletter
\newlabel{formal-spaces-section-completion}{{0AIX}{source}{}{sourcepointer-33}{}}
\newlabel{formal-spaces-section-completion@cref}{{[originalsection][0][]0AIX}{[1][0][]source}{}{}{}}
\makeatother
\crefname{originalproposition}{original proposition}{original propositions}
\Crefname{originalproposition}{Original proposition}{Original propositions}
\makeatletter
\newlabel{proposition-spaces-diagonal-representable-noetherian}{{0GE7}{source}{}{sourcepointer-34}{}}
\newlabel{proposition-spaces-diagonal-representable-noetherian@cref}{{[originalproposition][0][]0GE7}{[1][0][]source}{}{}{}}
\makeatother
\crefname{originallemma}{original lemma}{original lemmas}
\Crefname{originallemma}{Original lemma}{Original lemmas}
\makeatletter
\newlabel{restricted-lemma-base-change-formal-modification}{{0GDN}{source}{}{sourcepointer-35}{}}
\newlabel{restricted-lemma-base-change-formal-modification@cref}{{[originallemma][0][]0GDN}{[1][0][]source}{}{}{}}
\makeatother
\crefname{originaltheorem}{original theorem}{original theorems}
\Crefname{originaltheorem}{Original theorem}{Original theorems}
\makeatletter
\newlabel{restricted-theorem-dilatations}{{0GDU}{source}{}{sourcepointer-36}{}}
\newlabel{restricted-theorem-dilatations@cref}{{[originaltheorem][0][]0GDU}{[1][0][]source}{}{}{}}
\makeatother
\title{Limit preservation of the formal-contraction gluing functor}
\author{The Stacks Project Authors}
\date{stacks-project:artin.tex:local-snapshot-20261011; GNU FDL 1.2 or later}
\begin{document}
\maketitle
\noindent\textit{Editorial note.} The complete530-line author proof is retained; the task concerns the exact contraction functor in Situationcontractions, not all formal completion functors.\par\medskip
\noindent\textit{Editorial note.} The separately prepared0GI9 establishes one smooth neighbourhood from versality; it neither states nor supplies the present descent of completion-valued data.\par\medskip
\noindent\textit{Allowed standard prerequisites (editorial).}\begin{itemize}
\item Explicit prerequisite boundary. Noetherian formal algebraic spaces, formal modifications, inverse limits of finitely presented morphisms, proper coherent cohomology and formal functions, Artin–Rees and adic completion. The precise contraction functor and its local solution/equality lemmas are retained.
\end{itemize}\medskip
\ifdefined\hypertarget\hypertarget{sourcepointer-1}{}\fi
\noindent\textit{Source reference.} Original proposition 0GI3: focused-seeds/AGP/sources/stacks-snapshot-20261011/tags.txt, source lines 20603--20603. Stable Stacks tag for an expressly allowed foundational or notational input; the counted construction remains the two concrete core obligations. The retained local core is not replaced by this pointer. See standard\_prerequisites.\par\medskip
\ifdefined\hypertarget\hypertarget{sourcepointer-2}{}\fi
\noindent\textit{Source reference.} Original section 0GE1: focused-seeds/AGP/sources/stacks-snapshot-20261011/tags.txt, source lines 20461--20461. Stable Stacks tag for an expressly allowed foundational or notational input; the counted construction remains the two concrete core obligations. The retained local core is not replaced by this pointer. See standard\_prerequisites.\par\medskip
\ifdefined\hypertarget\hypertarget{sourcepointer-3}{}\fi
\noindent\textit{Source reference.} Original section 05N0: focused-seeds/AGP/sources/stacks-snapshot-20261011/tags.txt, source lines 7289--7289. Stable Stacks tag for an expressly allowed foundational or notational input; the counted construction remains the two concrete core obligations. The retained local core is not replaced by this pointer. See standard\_prerequisites.\par\medskip
\ifdefined\hypertarget\hypertarget{sourcepointer-4}{}\fi
\noindent\textit{Source reference.} Original lemma 0GE3: focused-seeds/AGP/sources/stacks-snapshot-20261011/tags.txt, source lines 20463--20463. Stable Stacks tag for an expressly allowed foundational or notational input; the counted construction remains the two concrete core obligations. The retained local core is not replaced by this pointer. See standard\_prerequisites.\par\medskip
\ifdefined\hypertarget\hypertarget{sourcepointer-5}{}\fi
\noindent\textit{Source reference.} Original proposition 05ER: focused-seeds/AGP/sources/stacks-snapshot-20261011/tags.txt, source lines 7000--7000. Stable Stacks tag for an expressly allowed foundational or notational input; the counted construction remains the two concrete core obligations. The retained local core is not replaced by this pointer. See standard\_prerequisites.\par\medskip
\ifdefined\hypertarget\hypertarget{sourcepointer-6}{}\fi
\noindent\textit{Source reference.} Original lemma 08AS: focused-seeds/AGP/sources/stacks-snapshot-20261011/tags.txt, source lines 10542--10542. Stable Stacks tag for an expressly allowed foundational or notational input; the counted construction remains the two concrete core obligations. The retained local core is not replaced by this pointer. See standard\_prerequisites.\par\medskip
\ifdefined\hypertarget\hypertarget{sourcepointer-7}{}\fi
\noindent\textit{Source reference.} Original lemma 073K: focused-seeds/AGP/sources/stacks-snapshot-20261011/tags.txt, source lines 9064--9064. Stable Stacks tag for an expressly allowed foundational or notational input; the counted construction remains the two concrete core obligations. The retained local core is not replaced by this pointer. See standard\_prerequisites.\par\medskip
\ifdefined\hypertarget\hypertarget{sourcepointer-8}{}\fi
\noindent\textit{Source reference.} Original theorem 08AZ: focused-seeds/AGP/sources/stacks-snapshot-20261011/tags.txt, source lines 10549--10549. Stable Stacks tag for an expressly allowed foundational or notational input; the counted construction remains the two concrete core obligations. The retained local core is not replaced by this pointer. See standard\_prerequisites.\par\medskip
\ifdefined\hypertarget\hypertarget{sourcepointer-9}{}\fi
\noindent\textit{Source reference.} Original lemma 0GH5: focused-seeds/AGP/sources/stacks-snapshot-20261011/tags.txt, source lines 20570--20570. Stable Stacks tag for an expressly allowed foundational or notational input; the counted construction remains the two concrete core obligations. The retained local core is not replaced by this pointer. See standard\_prerequisites.\par\medskip
\ifdefined\hypertarget\hypertarget{sourcepointer-10}{}\fi
\noindent\textit{Source reference.} Original lemma 0GH4: focused-seeds/AGP/sources/stacks-snapshot-20261011/tags.txt, source lines 20569--20569. Stable Stacks tag for an expressly allowed foundational or notational input; the counted construction remains the two concrete core obligations. The retained local core is not replaced by this pointer. See standard\_prerequisites.\par\medskip
\ifdefined\hypertarget\hypertarget{sourcepointer-11}{}\fi
\noindent\textit{Source reference.} Original lemma 0GH6: focused-seeds/AGP/sources/stacks-snapshot-20261011/tags.txt, source lines 20571--20571. Stable Stacks tag for an expressly allowed foundational or notational input; the counted construction remains the two concrete core obligations. The retained local core is not replaced by this pointer. See standard\_prerequisites.\par\medskip
\ifdefined\hypertarget\hypertarget{sourcepointer-12}{}\fi
\noindent\textit{Source reference.} Original lemma 0AQT: focused-seeds/AGP/sources/stacks-snapshot-20261011/tags.txt, source lines 13521--13521. Stable Stacks tag for an expressly allowed foundational or notational input; the counted construction remains the two concrete core obligations. The retained local core is not replaced by this pointer. See standard\_prerequisites.\par\medskip
\ifdefined\hypertarget\hypertarget{sourcepointer-13}{}\fi
\noindent\textit{Source reference.} Original lemma 0GD6: focused-seeds/AGP/sources/stacks-snapshot-20261011/tags.txt, source lines 20431--20431. Stable Stacks tag for an expressly allowed foundational or notational input; the counted construction remains the two concrete core obligations. The retained local core is not replaced by this pointer. See standard\_prerequisites.\par\medskip
\ifdefined\hypertarget\hypertarget{sourcepointer-14}{}\fi
\noindent\textit{Source reference.} Original section 01YV: focused-seeds/AGP/sources/stacks-snapshot-20261011/tags.txt, source lines 2570--2570. Stable Stacks tag for an expressly allowed foundational or notational input; the counted construction remains the two concrete core obligations. The retained local core is not replaced by this pointer. See standard\_prerequisites.\par\medskip
\ifdefined\hypertarget\hypertarget{sourcepointer-15}{}\fi
\noindent\textit{Source reference.} Original lemma 0CUF: focused-seeds/AGP/sources/stacks-snapshot-20261011/tags.txt, source lines 16100--16100. Stable Stacks tag for an expressly allowed foundational or notational input; the counted construction remains the two concrete core obligations. The retained local core is not replaced by this pointer. See standard\_prerequisites.\par\medskip
\ifdefined\hypertarget\hypertarget{sourcepointer-16}{}\fi
\noindent\textit{Source reference.} Original lemma 0A2P: focused-seeds/AGP/sources/stacks-snapshot-20261011/tags.txt, source lines 12712--12712. Stable Stacks tag for an expressly allowed foundational or notational input; the counted construction remains the two concrete core obligations. The retained local core is not replaced by this pointer. See standard\_prerequisites.\par\medskip
\ifdefined\hypertarget\hypertarget{sourcepointer-17}{}\fi
\noindent\textit{Source reference.} Original proposition 01ZC: focused-seeds/AGP/sources/stacks-snapshot-20261011/tags.txt, source lines 2587--2587. Stable Stacks tag for an expressly allowed foundational or notational input; the counted construction remains the two concrete core obligations. The retained local core is not replaced by this pointer. See standard\_prerequisites.\par\medskip
\ifdefined\hypertarget\hypertarget{sourcepointer-18}{}\fi
\noindent\textit{Source reference.} Original lemma 0GHZ: focused-seeds/AGP/sources/stacks-snapshot-20261011/tags.txt, source lines 20599--20599. Stable Stacks tag for an expressly allowed foundational or notational input; the counted construction remains the two concrete core obligations. The retained local core is not replaced by this pointer. See standard\_prerequisites.\par\medskip
\ifdefined\hypertarget\hypertarget{sourcepointer-19}{}\fi
\noindent\textit{Source reference.} Original lemma 049N: focused-seeds/AGP/sources/stacks-snapshot-20261011/tags.txt, source lines 5586--5586. Stable Stacks tag for an expressly allowed foundational or notational input; the counted construction remains the two concrete core obligations. The retained local core is not replaced by this pointer. See standard\_prerequisites.\par\medskip
\ifdefined\hypertarget\hypertarget{sourcepointer-20}{}\fi
\noindent\textit{Source reference.} Original lemma 01Z4: focused-seeds/AGP/sources/stacks-snapshot-20261011/tags.txt, source lines 2579--2579. Stable Stacks tag for an expressly allowed foundational or notational input; the counted construction remains the two concrete core obligations. The retained local core is not replaced by this pointer. See standard\_prerequisites.\par\medskip
\ifdefined\hypertarget\hypertarget{sourcepointer-21}{}\fi
\noindent\textit{Source reference.} Original theorem 0ARB: focused-seeds/AGP/sources/stacks-snapshot-20261011/tags.txt, source lines 13539--13539. Stable Stacks tag for an expressly allowed foundational or notational input; the counted construction remains the two concrete core obligations. The retained local core is not replaced by this pointer. See standard\_prerequisites.\par\medskip
\ifdefined\hypertarget\hypertarget{sourcepointer-22}{}\fi
\noindent\textit{Source reference.} Original lemma 0GI4: focused-seeds/AGP/sources/stacks-snapshot-20261011/tags.txt, source lines 20604--20604. Stable Stacks tag for an expressly allowed foundational or notational input; the counted construction remains the two concrete core obligations. The retained local core is not replaced by this pointer. See standard\_prerequisites.\par\medskip
\ifdefined\hypertarget\hypertarget{sourcepointer-23}{}\fi
\noindent\textit{Source reference.} Original lemma 0ARW: focused-seeds/AGP/sources/stacks-snapshot-20261011/tags.txt, source lines 13559--13559. Stable Stacks tag for an expressly allowed foundational or notational input; the counted construction remains the two concrete core obligations. The retained local core is not replaced by this pointer. See standard\_prerequisites.\par\medskip
\ifdefined\hypertarget\hypertarget{sourcepointer-24}{}\fi
\noindent\textit{Source reference.} Original proposition 03XX: focused-seeds/AGP/sources/stacks-snapshot-20261011/tags.txt, source lines 5164--5164. Stable Stacks tag for an expressly allowed foundational or notational input; the counted construction remains the two concrete core obligations. The retained local core is not replaced by this pointer. See standard\_prerequisites.\par\medskip
\ifdefined\hypertarget\hypertarget{sourcepointer-25}{}\fi
\noindent\textit{Source reference.} Original lemma 01ZM: focused-seeds/AGP/sources/stacks-snapshot-20261011/tags.txt, source lines 2599--2599. Stable Stacks tag for an expressly allowed foundational or notational input; the counted construction remains the two concrete core obligations. The retained local core is not replaced by this pointer. See standard\_prerequisites.\par\medskip
\ifdefined\hypertarget\hypertarget{sourcepointer-26}{}\fi
\noindent\textit{Source reference.} Original lemma 081F: focused-seeds/AGP/sources/stacks-snapshot-20261011/tags.txt, source lines 10215--10215. Stable Stacks tag for an expressly allowed foundational or notational input; the counted construction remains the two concrete core obligations. The retained local core is not replaced by this pointer. See standard\_prerequisites.\par\medskip
\ifdefined\hypertarget\hypertarget{sourcepointer-27}{}\fi
\noindent\textit{Source reference.} Original lemma 081E: focused-seeds/AGP/sources/stacks-snapshot-20261011/tags.txt, source lines 10214--10214. Stable Stacks tag for an expressly allowed foundational or notational input; the counted construction remains the two concrete core obligations. The retained local core is not replaced by this pointer. See standard\_prerequisites.\par\medskip
\ifdefined\hypertarget\hypertarget{sourcepointer-28}{}\fi
\noindent\textit{Source reference.} Original lemma 01Z0: focused-seeds/AGP/sources/stacks-snapshot-20261011/tags.txt, source lines 2575--2575. Stable Stacks tag for an expressly allowed foundational or notational input; the counted construction remains the two concrete core obligations. The retained local core is not replaced by this pointer. See standard\_prerequisites.\par\medskip
\ifdefined\hypertarget\hypertarget{sourcepointer-29}{}\fi
\noindent\textit{Source reference.} Original lemma 00L6: focused-seeds/AGP/sources/stacks-snapshot-20261011/tags.txt, source lines 776--776. Stable Stacks tag for an expressly allowed foundational or notational input; the counted construction remains the two concrete core obligations. The retained local core is not replaced by this pointer. See standard\_prerequisites.\par\medskip
\ifdefined\hypertarget\hypertarget{sourcepointer-30}{}\fi
\noindent\textit{Source reference.} Original section 0BNH: focused-seeds/AGP/sources/stacks-snapshot-20261011/tags.txt, source lines 14665--14665. Stable Stacks tag for an expressly allowed foundational or notational input; the counted construction remains the two concrete core obligations. The retained local core is not replaced by this pointer. See standard\_prerequisites.\par\medskip
\ifdefined\hypertarget\hypertarget{sourcepointer-31}{}\fi
\noindent\textit{Source reference.} Original lemma 0GC6: focused-seeds/AGP/sources/stacks-snapshot-20261011/tags.txt, source lines 20396--20396. Stable Stacks tag for an expressly allowed foundational or notational input; the counted construction remains the two concrete core obligations. The retained local core is not replaced by this pointer. See standard\_prerequisites.\par\medskip
\ifdefined\hypertarget\hypertarget{sourcepointer-32}{}\fi
\noindent\textit{Source reference.} Original definition 0GDK: focused-seeds/AGP/sources/stacks-snapshot-20261011/tags.txt, source lines 20445--20445. Stable Stacks tag for an expressly allowed foundational or notational input; the counted construction remains the two concrete core obligations. The retained local core is not replaced by this pointer. See standard\_prerequisites.\par\medskip
\ifdefined\hypertarget\hypertarget{sourcepointer-33}{}\fi
\noindent\textit{Source reference.} Original section 0AIX: focused-seeds/AGP/sources/stacks-snapshot-20261011/tags.txt, source lines 13280--13280. Stable Stacks tag for an expressly allowed foundational or notational input; the counted construction remains the two concrete core obligations. The retained local core is not replaced by this pointer. See standard\_prerequisites.\par\medskip
\ifdefined\hypertarget\hypertarget{sourcepointer-34}{}\fi
\noindent\textit{Source reference.} Original proposition 0GE7: focused-seeds/AGP/sources/stacks-snapshot-20261011/tags.txt, source lines 20467--20467. Stable Stacks tag for an expressly allowed foundational or notational input; the counted construction remains the two concrete core obligations. The retained local core is not replaced by this pointer. See standard\_prerequisites.\par\medskip
\ifdefined\hypertarget\hypertarget{sourcepointer-35}{}\fi
\noindent\textit{Source reference.} Original lemma 0GDN: focused-seeds/AGP/sources/stacks-snapshot-20261011/tags.txt, source lines 20448--20448. Stable Stacks tag for an expressly allowed foundational or notational input; the counted construction remains the two concrete core obligations. The retained local core is not replaced by this pointer. See standard\_prerequisites.\par\medskip
\ifdefined\hypertarget\hypertarget{sourcepointer-36}{}\fi
\noindent\textit{Source reference.} Original theorem 0GDU: focused-seeds/AGP/sources/stacks-snapshot-20261011/tags.txt, source lines 20454--20454. Stable Stacks tag for an expressly allowed foundational or notational input; the counted construction remains the two concrete core obligations. The retained local core is not replaced by this pointer. See standard\_prerequisites.\par\medskip
\begin{lemma}
\label{lemma-limit-preserving}
In Situation \ref{situation-contractions} the functor $F$ is limit preserving:
for any directed limit $V = \lim V_\lambda$ of Noetherian affine schemes
over $S$ we have $F(V) = \colim F(V_\lambda)$.
\end{lemma}

\begin{situation}
\label{situation-contractions}
Let $S$ be a locally Noetherian scheme. Let $X'$ be an algebraic space
locally of finite type over $S$. Let $T' \subset |X'|$ be a closed
subset. Let $U' \subset X'$ be the open subspace with
$|U'| = |X'| \setminus T'$. Let $W$ be a locally Noetherian
formal algebraic space over $S$ with $W_{red}$ locally of finite type
over $S$. Finally, we let
$$
g : X'_{/T'} \longrightarrow W
$$
be a formal modification, see
Algebraization of Formal Spaces, Definition
\ref{restricted-definition-formal-modification}.
Recall that $X'_{/T'}$ denotes the formal completion of $X'$ along
$T'$, see Formal Spaces, Section \ref{formal-spaces-section-completion}.
\end{situation}

\noindent
In the situation above our goal is to prove that there exists a
proper morphism $f : X' \to X$ of algebraic spaces over $S$,
a closed subset $T \subset |X|$, and an isomorphism
$a : X_{/T} \to W$ of formal algebraic spaces such that
\begin{enumerate}
\item $T'$ is the inverse image of $T$ by $|f| : |X'| \to |X|$,
\item $f : X' \to X$ maps $U'$ isomorphically to
an open subspace $U$ of $X$, and
\item $g = a \circ f_{/T}$ where $f_{/T} : X'_{/T'} \to X_{/T}$
is the induced morphism.
\end{enumerate}
Let us say that $(f : X' \to X, T, a)$ is a {\it solution}.

\medskip\noindent
We will follow Artin's strategy by constructing a functor $F$ on
the category of locally Noetherian schemes over $S$, showing that $F$ is an
algebraic space using
Proposition \ref{proposition-spaces-diagonal-representable-noetherian},
and proving that setting $X = F$ works.

\begin{remark}
\label{remark-G-rings}
In particular, we cannot prove that the desired result is true for
every Situation \ref{situation-contractions} because we will need to
assume the local rings of $S$ are G-rings. If you can prove the
result in general or if you have a counter example, please let
us know at
\href{mailto:stacks.project@gmail.com}{stacks.project@gmail.com}.
\end{remark}

\noindent
In Situation \ref{situation-contractions}
let $V$ be a locally Noetherian scheme over $S$.
The value of our functor $F$ on $V$ will be all triples
$$
(Z, u' : V \setminus Z \to U', \hat x : V_{/Z} \to W)
$$
satisfying the following conditions
\begin{enumerate}
\item $Z \subset V$ is a closed subset,
\item $u' : V \setminus Z \to U'$ is a morphism over $S$,
\item $\hat x : V_{/Z} \to W$ is an adic morphism of formal algebraic
spaces over $S$,
\item $u'$ and $\hat x$ are compatible (see below).
\end{enumerate}
The compatibility condition is the following: pulling back the
formal modification $g$ we obtain a formal modification
$$
X'_{/T'} \times_{g, W, \hat x} V_{/Z} \longrightarrow V_{/Z}
$$
See Algebraization of Formal Spaces, Lemma
\ref{restricted-lemma-base-change-formal-modification}.
By the main theorem on dilatations
(Algebraization of Formal Spaces, Theorem
\ref{restricted-theorem-dilatations}), there is a unique proper
morphism $V' \to V$ of algebraic spaces which is an isomorphism over
$V \setminus Z$ such that $V'_{/Z} \to V_{/Z}$ is isomorphic to the
displayed arrow. In other words, for some morphism
$\hat x' : V'_{/Z} \to X'_{/T'}$ we have a cartesian diagram
$$
\xymatrix{
V'_{/Z} \ar[r] \ar[d]_{\hat x'} & V_{/Z} \ar[d]^{\hat x} \\
X'_{/T'} \ar[r]^g & W
}
$$
of formal algebraic spaces. We will think
of $V \setminus Z$ as an open subspace of $V'$ without further mention.
The compatibility condition is that there should be a
morphism $x' : V' \to X'$ restricting to $u'$ and $\hat x$
over $V \setminus Z \subset V'$ and $V'_{/Z}$.
In other words, such that the diagram
$$
\xymatrix{
V \setminus Z \ar[r] \ar[d]_{u'} &
V' \ar[d]^{x'} &
V'_{/Z} \ar[l] \ar[d]^{\hat x'} \ar[r] &
V_{/Z} \ar[d]^{\hat x} \\
U' \ar[r] &
X' &
X'_{/T'} \ar[r]^g \ar[l] &
W
}
$$
is commutative. Observe that by Algebraization of Formal Spaces,

if it exists. We will indicate this situation by saying
``{\it $V' \to V$, $\hat x'$, and $x'$ witness the compatibility
between $u'$ and $\hat x$}''.

\begin{remark}
\label{remark-how-to-think-compatibility}
In Situation \ref{situation-contractions} let $V$ be a locally Noetherian
scheme over $S$. Let $(Z, u', \hat x)$ be a triple satisfying (1), (2), and
(3) above. We want to explain a way to think about the compatibility
condition (4). It will not be mathematically precise as we are going use
a fictitious category $\textit{An}_S$ of analytic spaces over $S$
and a fictitious analytification functor
$$
\left\{
\begin{matrix}
\text{locally Noetherian formal} \\
\text{algebraic spaces over }S
\end{matrix}
\right\}
\longrightarrow
\textit{An}_S,
\quad\quad
Y \longmapsto Y^{an}
$$
For example if $Y = \text{Spf}(k[[t]])$ over $S = \Spec(k)$, then $Y^{an}$
should be thought of as an open unit disc. If $Y = \Spec(k)$, then $Y^{an}$
is a single point. The category $\textit{An}_S$ should have open and
closed immersions and we should be able to take the open complement
of a closed. Given $Y$ the morphism $Y_{red} \to Y$ should induces a
closed immersion $Y_{red}^{an} \to Y^{an}$. We set
$Y^{rig} = Y^{an} \setminus Y_{red}^{an}$ equal to its open complement.
If $Y$ is an algebraic space and if $Z \subset Y$ is closed, then
the morphism $Y_{/Z} \to Y$ should induce an open immersion
$Y_{/Z}^{an} \to Y^{an}$ which in turn should induce an open immersion
$$
can : (Y_{/Z})^{rig} \longrightarrow (Y \setminus Z)^{an}
$$
Also, given a formal modification $g : Y' \to Y$ of locally Noetherian formal
algebraic spaces, the induced morphism $g^{rig} : (Y')^{rig} \to Y^{rig}$
should be an isomorphism. Given $\text{An}_S$ and the analytification
functor, we can consider the requirement that
$$
\xymatrix{
(V_{/Z})^{rig} \ar[rr]_{can} \ar[d]_{(g^{rig})^{-1} \circ \hat x^{an}} & &
(V \setminus Z)^{an} \ar[d]^{(u')^{an}} \\
(X'_{/T'})^{rig} \ar[rr]^{can} & & (X' \setminus T')^{an}
}
$$
commutes. This makes sense as $g^{rig} : (X'_{T'})^{rig} \to W^{rig}$
is an isomorphism and $U' = X' \setminus T'$. Finally, under some assumptions
of faithfulness of the analytification functor, this requirement will
be equivalent to the compatibility condition formulated above.
We hope this will motivate the reader to think of the compatibility
of $u'$ and $\hat x$ as the requirement that some maps be equal,
rather than asking for the existence of a certain commutative diagram.
\end{remark}

\begin{lemma}
\label{lemma-functor}
In Situation \ref{situation-contractions} the rule $F$ that sends
a locally Noetherian scheme $V$ over $S$ to the set of triples
$(Z, u', \hat x)$ satisfying the compatibility condition and which sends a
a morphism $\varphi : V_2 \to V_1$ of locally Noetherian schemes over $S$
to the map
$$
F(\varphi) : F(V_1) \longrightarrow F(V_2)
$$
sending an element $(Z_1, u'_1, \hat x_1)$ of $F(V_1)$ to
$(Z_2, u'_2, \hat x_2)$ in $F(V_2)$ given by
\begin{enumerate}
\item $Z_2 \subset V_2$ is the inverse image of $Z_1$ by $\varphi$,
\item $u'_2$ is the composition of $u'_1$ and
$\varphi|_{V_2 \setminus Z_2} : V_2 \setminus Z_2 \to V_1 \setminus Z_1$,
\item $\hat x_2$ is the composition of $\hat x_1$ and
$\varphi_{/Z_2} : V_{2, /Z_2} \to V_{1, /Z_1}$
\end{enumerate}
is a contravariant functor.
\end{lemma}

\begin{proof}
To see the compatibility condition between $u'_2$ and $\hat x_2$, let
$V'_1 \to V_1$, $\hat x'_1$, and $x'_1$ witness the compatibility between
$u'_1$ and $\hat x_1$. Set $V'_2 = V_2 \times_{V_1} V'_1$, set
$\hat x'_2$ equal to the composition of $\hat x'_1$ and
$V'_{2, /Z_2} \to V'_{1, /Z_1}$, and set $x'_2$
equal to the composition of $x'_1$ and $V'_2 \to V'_1$.
Then $V'_2 \to V_2$, $\hat x'_2$, and $x'_2$ witness the compatibility between
$u'_2$ and $\hat x_2$. We omit the detailed verification.
\end{proof}

\begin{lemma}
\label{lemma-solution}
In Situation \ref{situation-contractions} if there exists a solution
$(f : X' \to X, T, a)$ then there is a functorial bijection
$F(V) = \Mor_S(V, X)$ on the category of
locally Noetherian schemes $V$ over $S$.
\end{lemma}

\begin{proof}
Let $V$ be a locally Noetherian scheme over $S$.
Let $x : V \to X$ be a morphism over $S$.
Then we get an element $(Z, u', \hat x)$ in $F(V)$ as follows
\begin{enumerate}
\item $Z \subset V$ is the inverse image of $T$ by $x$,
\item $u' : V \setminus Z \to U' = U$ is the restriction of
$x$ to $V \setminus Z$,
\item $\hat x : V_{/Z} \to W$ is the composition of
$x_{/Z} : V_{/Z} \to X_{/T}$ with the isomorphism $a : X_{/T} \to W$.
\end{enumerate}
This triple satisfies the compatibility condition because we
can take $V' = V \times_{x, X} X'$, we can take $\hat x'$
the completion of the projection $x' : V' \to X'$.

\medskip\noindent
Conversely, suppose given an element $(Z, u', \hat x)$ of $F(V)$. We claim
there is a unique morphism $x : V \to X$ compatible with $u'$ and $\hat x$.
Namely, let $V' \to V$, $\hat x'$, and $x'$ witness the
compatibility between $u'$ and $\hat x$. Then
Algebraization of Formal Spaces, Proposition
\ref{restricted-proposition-glue-modification}
is exactly the result we need to find
a unique morphism $x : V \to X$ agreeing with
$\hat x$ over $V_{/Z}$ and with $x'$ over $V'$ (and a fortiori
agreeing with $u'$ over $V \setminus Z$).

\medskip\noindent
We omit the verification that the two constructions above define inverse
bijections between their respective domains.
\end{proof}

\begin{lemma}
\label{lemma-functor-is-solution}
In Situation \ref{situation-contractions} if there exists an
algebraic space $X$ locally of finite type over $S$ and a
functorial bijection $F(V) = \Mor_S(V, X)$ on the category of
locally Noetherian schemes $V$ over $S$, then $X$ is a solution.
\end{lemma}

\begin{proof}
We have to construct a proper morphism $f : X' \to X$, a closed subset
$T \subset |X|$, and an isomorphism $a : X_{/T} \to W$ with properties
(1), (2), (3) listed just below Situation \ref{situation-contractions}.

\medskip\noindent
The discussion in this proof is a bit pedantic because we want
to carefully match the underlying categories. In this paragraph
we explain how the adventurous reader can proceed less timidly.
Namely, the reader may extend our definition of the functor $F$
to all locally Noetherian algebraic spaces over $S$.
Doing so the reader may then conclude that $F$ and $X$ agree as
functors on the category of these algebraic spaces, i.e.,
$X$ represents $F$. Then one considers the universal object
$(T, u', \hat x)$ in $F(X)$. Then the reader will find that
for the triple $X'' \to X$, $\hat x'$, $x'$ witnessing the compatibility
between $u'$ and $\hat x$ the morphism $x' : X'' \to X'$ is an isomorphism
and this will produce $f : X' \to X$ by inverting $x'$. Finally, we already
have $T \subset |X|$ and the reader may show that $\hat x$ is an isomorphism
which can served as the last ingredient namely $a$.

\medskip\noindent
Denote $h_X(-) = \Mor_S(-, X)$ the functor of
points of $X$ restricted to the category
$(\textit{Noetherian}/S)_\etale$ of Section \ref{section-noetherian}.
By Limits of Spaces, Remark \ref{spaces-limits-remark-limit-preserving}
the algebraic spaces $X$ and $X'$ are limit preserving. Hence so
are the restrictions $h_X$ and $h_{X'}$.
To construct $f$ it therefore suffices to construct a
transformation $h_{X'} \to h_X = F$, see Lemma \ref{lemma-canonical-extension}.
Thus let $V \to S$ be an object of $(\textit{Noetherian}/S)_\etale$
and let $\tilde x : V \to X'$ be in $h_{X'}(V)$.
Then we get an element $(Z, u', \hat x)$ in $F(V)$ as follows
\begin{enumerate}
\item $Z \subset V$ is the inverse image of $T'$ by $\tilde x$,
\item $u' : V \setminus Z \to U'$ is the restriction of
$\tilde x$ to $V \setminus Z$,
\item $\hat x : V_{/Z} \to W$ is the composition of
$x_{/Z} : V_{/Z} \to X'_{/T'}$ with $g : X'_{/T'} \to W$.
\end{enumerate}
This triple satisfies the compatibility condition: first we always
obtain $V' \to V$ and $\hat x' : V'_{/Z'} \to X'_{/T'}$ for free
(see discussion preceding Lemma \ref{lemma-functor}).
Then we just define $x' : V' \to X'$ to be the composition
of $V' \to V$ and the morphism $\tilde x : V \to X'$.
We omit the verification that this works.

\medskip\noindent
If $\xi : V \to X$ is an \'etale morphism where $V$ is a scheme, then
we obtain $\xi = (Z, u', \hat x) \in F(V) = h_X(V) = X(V)$.
Of course, if $\varphi : V' \to V$ is a further \'etale morphism of schemes,
then $(Z, u', \hat x)$ pulled back to $F(V')$ corresponds to
$\xi \circ \varphi$.
The closed subset $T \subset |X|$ is just defined as the closed
subset such that $\xi : V \to X$ for $\xi = (Z, u', \hat x)$
pulls $T$ back to $Z$

\medskip\noindent
Consider Noetherian schemes $V$ over $S$ and a morphism
$\xi : V \to X$ corresponding to $(Z, u', \hat x)$ as above.
Then we see that $\xi(V)$ is set theoretically contained in $T$
if and only if $V = Z$ (as topological spaces). Hence we see that
$X_{/T}$ agrees with $W$ as a functor. This produces the isomorphism
$a : X_{/T} \to W$. (We've omitted a small detail here which is
that for the locally Noetherian formal algebraic spaces $X_{/T}$ and
$W$ it suffices to check one gets an isomorphism after evaluating
on locally Noetherian schemes over $S$.)

\medskip\noindent
We omit the proof of conditions (1), (2), and (3).
\end{proof}

\begin{remark}
\label{remark-diagonal}
In Situation \ref{situation-contractions}.
Let $V$ be a locally Noetherian scheme over $S$.
Let $(Z_i, u'_i, \hat x_i) \in F(V)$ for $i = 1, 2$. Let $V'_i \to V$,
$\hat x'_i$ and $x'_i$ witness the compatibility between $u'_i$ and
$\hat x_i$ for $i = 1, 2$.

\medskip\noindent
Set $V' = V'_1 \times_V V'_2$. Let $E' \to V'$ denote the equalizer
of the morphisms
$$
V' \to V'_1 \xrightarrow{x'_1} X'
\quad\text{and}\quad
V' \to V'_2 \xrightarrow{x'_2} X'
$$
Set $Z = Z_1 \cap Z_2$. Let $E_W \to V_{/Z}$
be the equalizer of the morphisms
$$
V_{/Z} \to V_{/Z_1} \xrightarrow{\hat x_1} W
\quad\text{and}\quad
V_{/Z} \to V_{/Z_2} \xrightarrow{\hat x_2} W
$$
Observe that $E' \to V$ is separated and locally of finite type
and that $E_W$ is a locally Noetherian formal algebraic space
separated over $V$.
The compatibilities between the various morphisms involved show that
\begin{enumerate}
\item $\Im(E' \to V) \cap (Z_1 \cup Z_2)$
is contained in $Z = Z_1 \cap Z_2$, 
\item the morphism $E' \times_V (V \setminus Z) \to V \setminus Z$
is a monomorphism and is equal to the equalizer of the restrictions
of $u'_1$ and $u'_2$ to $V \setminus (Z_1 \cup Z_2)$,
\item the morphism $E'_{/Z} \to V_{/Z}$ factors through $E_W$
and the diagram
$$
\xymatrix{
E'_{/Z} \ar[r] \ar[d] & X'_{/T'} \ar[d]^g \\
E_W \ar[r] & W
}
$$
is cartesian. In particular, the morphism $E'_{/Z} \to E_W$
is a formal modification as the base change of $g$,
\item $E'$, $(E' \to V)^{-1}Z$, and $E'_{/Z} \to E_W$
is a triple as in Situation \ref{situation-contractions}
with base scheme the locally Noetherian scheme $V$,
\item given a morphism $\varphi : A \to V$
of locally Noetherian schemes, the following are equivalent
\begin{enumerate}
\item $(Z_1, u'_1, \hat x_1)$ and $(Z_2, u'_2, \hat x_2)$
restrict to the same element of $F(A)$,
\item $A \setminus \varphi^{-1}(Z) \to V \setminus Z$
factors through $E' \times_V (V \setminus Z)$
and $A_{/\varphi^{-1}(Z)} \to V_{/Z}$
factors through $E_W$.
\end{enumerate}
\end{enumerate}
We conclude, using
Lemmas \ref{lemma-solution} and \ref{lemma-functor-is-solution},
that if there is a solution $E \to V$
for the triple in (4), then $E$ represents
$F \times_{\Delta, F \times F} V$ on the category of
locally Noetherian schemes over $V$.
\end{remark}

\begin{lemma}
\label{lemma-closed-immersion}
In Situation \ref{situation-contractions} assume given a closed
subset $Z \subset S$ such that
\begin{enumerate}
\item the inverse image of $Z$ in $X'$ is $T'$,
\item $U' \to S \setminus Z$ is a closed immersion,
\item $W \to S_{/Z}$ is a closed immersion.
\end{enumerate}
Then there exists a solution $(f : X' \to X, T, a)$
and moreover $X \to S$ is a closed immersion.
\end{lemma}

\begin{proof}
Suppose we have a closed subscheme $X \subset S$ such that
$X \cap (S \setminus Z) = U'$ and $X_{/Z} = W$. Then
$X$ represents the functor $F$ (some details omitted) and hence 
is a solution. To find $X$ is clearly a local question on $S$.
In this way we reduce to the case discussed in the next paragraph.

\medskip\noindent
Assume $S = \Spec(A)$ is affine. Let $I \subset A$ be the
radical ideal cutting out $Z$.
Write $I = (f_1, \ldots, f_r)$. By assumption we are given
\begin{enumerate}
\item the closed immersion $U' \to S \setminus Z$ determines
ideals $J_i \subset A[1/f_i]$ such that $J_i$ and $J_j$
generate the same ideal in $A[1/f_if_j]$,
\item the closed immersion $W \to S_{/Z}$ is the map
$\text{Spf}(A^\wedge/J') \to \text{Spf}(A^\wedge)$ for some
ideal $J' \subset A^\wedge$ in the $I$-adic completion $A^\wedge$ of $A$.
\end{enumerate}
To finish the proof we need to find an ideal $J \subset A$
such that $J_i = J[1/f_i]$ and $J' = JA^\wedge$. By
More on Algebra, Proposition \ref{more-algebra-proposition-equivalence}
it suffices to show that $J_i$ and $J'$ generate the same ideal
in $A^\wedge[1/f_i]$ for all $i$.

\medskip\noindent
Recall that $A' = H^0(X', \mathcal{O})$ is a finite $A$-algebra
whose formation commutes with flat base change
(Cohomology of Spaces, Lemmas
\ref{spaces-cohomology-lemma-proper-over-affine-cohomology-finite} and
\ref{spaces-cohomology-lemma-flat-base-change-cohomology}). Denote
$J'' = \Ker(A \to A')$\footnote{Contrary to what the reader
may expect, the ideals $J$ and $J''$ won't agree in general.}.
We have $J_i = J''A[1/f_i]$ as follows
from base change to the spectrum of $A[1/f_i]$.
Observe that we have a commutative diagram
$$
\xymatrix{
X' \ar[d] &
X'_{/T'} \times_{S_{/Z}} \text{Spf}(A^\wedge) \ar[l] \ar[d] &
X'_{/T'} \times_W \text{Spf}(A^\wedge/J') \ar@{=}[l] \ar[d] \\
\Spec(A) &
\text{Spf}(A^\wedge) \ar[l] &
\text{Spf}(A^\wedge/J') \ar[l]
}
$$
The middle vertical arrow is the completion of the left vertical
arrow along the obvious closed subsets. By the theorem on formal
functions we have
$$
(A')^\wedge = \Gamma(X' \times_S \Spec(A^\wedge), \mathcal{O}) =
\lim H^0(X' \times_S \Spec(A/I^n), \mathcal{O})
$$
See Cohomology of Spaces, Theorem
\ref{spaces-cohomology-theorem-formal-functions}.
From the diagram we conclude that $J'$ maps to zero in $(A')^\wedge$.
Hence $J' \subset J'' A^\wedge$. Consider the arrows
$$
X'_{/T'} \to
\text{Spf}(A^\wedge/J''A^\wedge) \to
\text{Spf}(A^\wedge/J') = W
$$
We know the composition $g$ is a formal modification
(in particular rig-\'etale and rig-surjective) and the second
arrow is a closed immersion (in particular an adic monomorphism).
Hence $X'_{/T'} \to \text{Spf}(A^\wedge/J''A^\wedge)$ is
rig-surjective and rig-\'etale, see
Algebraization of Formal Spaces, Lemmas
\ref{restricted-lemma-rig-surjective-alternative-permanence} and
\ref{restricted-lemma-rig-etale-alternative-permanence}.
Applying Algebraization of Formal Spaces, Lemmas
\ref{restricted-lemma-rig-etale-descent} and
\ref{restricted-lemma-permanence-rig-surjective}
we conclude that $\text{Spf}(A^\wedge/J''A^\wedge) \to W$
is rig-\'etale and rig-surjective.
By Algebraization of Formal Spaces, Lemma
\ref{restricted-lemma-closed-immersion-rig-smooth-rig-surjective}
we conclude that $I^n J'' A^\wedge \subset J'$ for some $n > 0$.
It follows that $J'' A^\wedge[1/f_i] = J' A^\wedge[1/f_i]$ and
we deduce $J_i A^\wedge[1/f_i] = J' A^\wedge[1/f_i]$ for all
$i$ as desired.
\end{proof}

\begin{lemma}
\label{lemma-diagonal-contractions}
In Situation \ref{situation-contractions} assume $X' \to S$
and $W \to S$ are separated. Then the diagonal $\Delta : F \to F \times F$
is representable by closed immersions.
\end{lemma}

\begin{proof}
Combine Lemma \ref{lemma-closed-immersion}
with the discussion in Remark \ref{remark-diagonal}.
\end{proof}

\begin{lemma}
\label{lemma-sheaf}
In Situation \ref{situation-contractions} the functor
$F$ satisfies the sheaf property for all \'etale coverings
of locally Noetherian schemes over $S$.
\end{lemma}

\begin{proof}
Omitted. Hint: morphisms may be defined \'etale locally.
\end{proof}

\begin{proof}
This is an absurdly long proof. Much of it consists of standard
arguments on limits and \'etale localization. We urge the reader to skip ahead
to the last part of the proof where something interesting happens.

\medskip\noindent
Let $V = \lim_{\lambda \in \Lambda} V_i$ be a directed limit of schemes
over $S$ with $V$ and $V_\lambda$ Noetherian and with affine transition
morphisms. See Limits, Section \ref{limits-section-limits} for material on
limits of schemes. We will prove that $\colim F(V_\lambda) \to F(V)$
is bijective.

\medskip\noindent
Proof of injectivity: notation.
Let $\lambda \in \Lambda$ and
$\xi_{\lambda, 1}, \xi_{\lambda, 2} \in F(V_\lambda)$ 
be elements which restrict to the same element of $F(V)$.
Write $\xi_{\lambda, 1} = (Z_{\lambda, 1}, u'_{\lambda, 1},
\hat x_{\lambda, 1})$ and
$\xi_{\lambda, 2} = (Z_{\lambda, 2}, u'_{\lambda, 2}, \hat x_{\lambda, 2})$.

\medskip\noindent
Proof of injectivity: agreement of $Z_{\lambda, i}$.
Since $Z_{\lambda, 1}$ and $Z_{\lambda, 2}$ restrict to the same closed
subset of $V$, we may after increasing $i$ assume
$Z_{\lambda, 1} = Z_{\lambda, 2}$, see
Limits, Lemma \ref{limits-lemma-inverse-limit-top} and
Topology, Lemma \ref{topology-lemma-describe-limits}.
Let us denote the common value $Z_\lambda \subset V_\lambda$, for
$\mu \geq \lambda$
denote $Z_\mu \subset V_\mu$ the inverse image in $V_\mu$ and
and denote $Z$ the inverse image in $V$. We will use below that
$Z = \lim_{\mu \geq \lambda} Z_\mu$ as schemes if we view $Z$
and $Z_\mu$ as reduced closed subschemes.

\medskip\noindent
Proof of injectivity: agreement of $u'_{\lambda, i}$.
Since $U'$ is locally of finite type over $S$ and since
the restrictions of $u'_{\lambda, 1}$ and $u'_{\lambda, 2}$
to $V \setminus Z$ are the same, we may after increasing $\lambda$ assume
$u'_{\lambda, 1} = u'_{\lambda, 2}$, see Limits, Proposition
\ref{limits-proposition-characterize-locally-finite-presentation}.
Let us denote the common value $u'_\lambda$ and denote
$u'$ the restriction to $V \setminus Z$.

\medskip\noindent
Proof of injectivity: restatement.
At this point we have
$\xi_{\lambda, 1} = (Z_\lambda, u'_\lambda, \hat x_{\lambda, 1})$ and
$\xi_{\lambda, 2} = (Z_\lambda, u'_\lambda, \hat x_{\lambda, 2})$.
The main problem we face in this part of the proof
is to show that the morphisms
$\hat x_{\lambda, 1}$ and $\hat x_{\lambda, 2}$ become the same after
increasing $\lambda$.

\medskip\noindent
Proof of injectivity: agreement of $\hat x_{\lambda, i}|_{Z_\lambda}$.
Consider the morpisms
$\hat x_{\lambda, 1}|_{Z_\lambda}, \hat x_{\lambda, 2}|_{Z_\lambda} :
Z_\lambda \to W_{red}$.
These morphisms restrict to the same morphism $Z \to W_{red}$.
Since $W_{red}$ is a scheme locally of finite type over $S$
we see using Limits, Proposition
\ref{limits-proposition-characterize-locally-finite-presentation}
that after replacing $\lambda$ by a bigger index
we may assume
$\hat x_{\lambda, 1}|_{Z_\lambda} = \hat x_{\lambda, 2}|_{Z_\lambda} :
Z_\lambda \to W_{red}$.

\medskip\noindent
Proof of injectivity: end.
Next, we are going to apply the discussion in
Remark \ref{remark-diagonal} to $V_\lambda$ and the two elements
$\xi_{\lambda, 1}, \xi_{\lambda, 2} \in F(V_\lambda)$.
This gives us
\begin{enumerate}
\item $e_\lambda : E_\lambda' \to V_\lambda$
separated and locally of finite type,
\item $e_\lambda^{-1}(V_\lambda \setminus Z_\lambda) \to
V_\lambda \setminus Z_\lambda$ is an isomorphism,
\item a monomorphism $E_{W, \lambda} \to V_{\lambda, /Z_\lambda}$
which is the equalizer of $\hat x_{\lambda, 1}$ and $\hat x_{\lambda, 2}$,
\item a formal modification $E'_{\lambda, /Z_\lambda} \to E_{W, \lambda}$
\end{enumerate}
Assertion (2) holds by assertion (2) in Remark \ref{remark-diagonal}
and the preparatory work we did above getting
$u'_{\lambda, 1} = u'_{\lambda, 2} = u'_\lambda$.
Since $Z_\lambda = (V_{\lambda, /Z_\lambda})_{red}$ factors through
$E_{W, \lambda}$ because
$\hat x_{\lambda, 1}|_{Z_\lambda} = \hat x_{\lambda, 2}|_{Z_\lambda}$
we see from
Formal Spaces, Lemma \ref{formal-spaces-lemma-monomorphism-iso-over-red}
that $E_{W, \lambda} \to V_{\lambda, /Z_\lambda}$ is a closed immersion.
Then we see from assertion (4) in Remark \ref{remark-diagonal}
and Lemma \ref{lemma-closed-immersion} applied to the triple
$E_\lambda'$, $e_\lambda^{-1}(Z_\lambda)$,
$E'_{\lambda, /Z_\lambda} \to E_{W, \lambda}$ over $V_\lambda$ that
there exists a closed immersion $E_\lambda \to V_\lambda$
which is a solution for this triple.
Next we use assertion (5) in Remark \ref{remark-diagonal}
which combined with Lemma \ref{lemma-solution}
says that $E_\lambda$ is the ``equalizer'' of $\xi_{\lambda, 1}$
and $\xi_{\lambda, 2}$. In particular, we see that $V \to V_\lambda$
factors through $E_\lambda$. Then using Limits, Proposition
\ref{limits-proposition-characterize-locally-finite-presentation}
once more we find $\mu \geq \lambda$ such that $V_\mu \to V_\lambda$
factors through $E_\lambda$ and hence the pullbacks of
$\xi_{\lambda, 1}$ and $\xi_{\lambda, 2}$ to $V_\mu$ are the same
as desired.

\medskip\noindent
Proof of surjectivity: statement.
Let $\xi = (Z, u', \hat x)$ be an element of $F(V)$.
We have to find a $\lambda \in \Lambda$ and an element
$\xi_\lambda \in F(V_\lambda)$ restricting to $\xi$.

\medskip\noindent
Proof of surjectivity: the question is \'etale local.
By the unicity proved in the previous part of the proof and by the
sheaf property of $F$ in Lemma \ref{lemma-sheaf}, the problem
is local on $V$ in the \'etale topology. More precisely, let $v \in V$.
We claim it suffices to find an \'etale morphism
$(\tilde V, \tilde v) \to (V, v)$ and some
$\lambda$, some an \'etale morphism $\tilde V_\lambda \to V_\lambda$,
and some element $\tilde \xi_\lambda \in F(\tilde V_\lambda)$ such that
$\tilde V = \tilde V_\lambda \times_{V_\lambda} V$
and $\xi|_U = \tilde \xi_\lambda|_U$. We omit a detailed proof of
this claim\footnote{To prove this
one assembles a collection of the morphisms $\tilde V \to V$
into a finite \'etale covering and shows that the corresponding morphisms
$\tilde V_\lambda \to V_\lambda$ form an \'etale covering as well (after
increasing $\lambda$). Next one uses the injectivity to see that
the elements $\tilde \xi_\lambda$ glue (after increasing $\lambda$)
and one uses the sheaf property for $F$ to descend these
elements to an element of $F(V_\lambda)$.}.

\medskip\noindent
Proof of surjectivity: rephrasing the problem.
Recall that any \'etale morphism $(\tilde V, \tilde v) \to (V, v)$
with $\tilde V$ affine is the base change of an \'etale morphism
$\tilde V_\lambda \to V_\lambda$ with $\tilde V_\lambda$ affine
for some $\lambda$, see for example
Topologies, Lemma \ref{topologies-lemma-limit-fppf-topology}.
Given $\tilde V_\lambda$ we have
$\tilde V = \lim_{\mu \geq \lambda} \tilde V_\lambda \times_{V_\lambda} V_\mu$.
Hence given $(\tilde V, \tilde v) \to (V, v)$ \'etale with $\tilde V$ affine,
we may replace $(V, v)$ by $(\tilde V, \tilde v)$ and $\xi$ by
the restriction of $\xi$ to $\tilde V$.

\medskip\noindent
Proof of surjectivity: reduce to base being affine. In particular,
suppose $\tilde S \subset S$ is an affine open subscheme such
that $v \in V$ maps to a point of $\tilde S$. Then we may according
to the previous paragraph, replace $V$ by $\tilde V = \tilde S \times_S V$.
Of course, if we do this, it suffices to solve the problem
for the functor $F$ restricted to the category of locally Noetherian
schemes over $\tilde S$. This functor is of course the functor
associated to the whole situation base changed to $\tilde S$.
Thus we may and do assume $S = \Spec(R)$ is a Noetherian affine scheme
for the rest of the proof.

\medskip\noindent
Proof of surjectivity: easy case.
If $v \in V \setminus Z$, then we can take $\tilde V = V \setminus Z$.
This descends to an open subscheme $\tilde V_\lambda \subset V_\lambda$
for some $\lambda$ by Limits, Lemma
\ref{limits-lemma-descend-opens}.
Next, after increasing $\lambda$ we may assume
there is a morphism $u'_\lambda : \tilde V_\lambda \to U'$
restricting to $u'$. Taking
$\tilde \xi_\lambda = (\emptyset, u'_\lambda, \emptyset)$
gives the desired element of $F(\tilde V_\lambda)$.

\medskip\noindent
Proof of surjectivity: hard case and reduction to affine $W$.
The most difficult case comes from considering $v \in Z \subset V$.
We claim that we can reduce this to the case where $W$ is an
affine formal scheme; we urge the reader to skip this argument\footnote{Artin's
approach to the proof of this lemma is to work around this and
consequently he can avoid proving the injectivity first. Namely, Artin
consistently works with a finite affine \'etale coverings of all spaces
in sight keeping track of the maps between them during the proof.
In hindsight that might be preferable to what we do here.}.
Namely, we can choose an \'etale morphism
$\tilde W \to W$ where $\tilde W$ is an affine formal algebraic space
such that the image of $v$ by $\hat x : V_{/Z} \to W$ is
in the image of $\tilde W \to W$ (on reductions).
Then the morphisms
$$
p : \tilde W \times_{W, g} X'_{/T'} \longrightarrow X'_{/T'}
$$
and
$$
q : \tilde W \times_{W, \hat x} V_{/Z} \to V_{/Z}
$$
are \'etale morphisms of locally Noetherian formal algebraic spaces.
By (an easy case of) Algebraization of Formal Spaces, Theorem
\ref{restricted-theorem-dilatations-general}
there exists a morphism $\tilde X' \to X'$ of algebraic spaces
which is locally of finite type, is an isomorphism over $U'$, and
such that $\tilde X'_{/T'} \to X'_{/T'}$ is isomorphic to $p$.
By Algebraization of Formal Spaces, Lemma \ref{restricted-lemma-output-etale}
the morphism $\tilde X' \to X'$ is \'etale. Denote
$\tilde T' \subset |\tilde X'|$ the inverse image of $T'$.
Denote $\tilde U' \subset \tilde X'$ the complementary open subspace.
Denote $\tilde g' : \tilde X'_{/\tilde T'} \to \tilde W$
the formal modification which is the base change of $g$ by
$\tilde W \to W$. Then we see that
$$
\tilde X',\ \tilde T',\ \tilde U',\ \tilde W,
\ \tilde g : \tilde X'_{/\tilde T'} \to \tilde W
$$
is another example of Situation \ref{situation-contractions}.
Denote $\tilde F$ the functor constructed from this triple.
There is a transformation of functors
$$
\tilde F \longrightarrow F
$$
constructed using the morphisms $\tilde X' \to X'$ and
$\tilde W \to W$ in the obvious manner; details omitted.

\medskip\noindent
Proof of surjectivity: hard case and reduction to affine $W$, part 2.
By the same theorem as used above, there exists a morphism
$\tilde V \to V$ of algebraic spaces which is locally of finite type,
is an isomorphism over $V \setminus Z$ and such that
$\tilde V_{/Z} \to V_{/Z}$ is isomorphic to $q$.
Denote $\tilde Z \subset \tilde V$ the inverse image of $Z$.
By Algebraization of Formal Spaces, Lemmas
\ref{restricted-lemma-output-etale} and \ref{restricted-lemma-output-separated}
the morphism $\tilde V \to V$ is \'etale and separated.
In particular $\tilde V$ is a (locally Noetherian) scheme, see for example
Morphisms of Spaces, Proposition
\ref{spaces-morphisms-proposition-locally-quasi-finite-separated-over-scheme}.
We have the morphism $u'$ which we may view as a morphism
$$
\tilde u' : \tilde V \setminus \tilde Z \longrightarrow \tilde U'
$$
where $\tilde U' \subset \tilde X'$ is the open mapping isomorphically
to $U'$. We have a morphism
$$
\tilde {\hat x} :
\tilde V_{/\tilde Z} = \tilde W \times_{W, \hat x} V_{/Z}
\longrightarrow
\tilde W
$$
Namely, here we just use the projection. Thus we have the triple
$$
\tilde \xi = (\tilde Z, \tilde u', \tilde {\hat x}) \in \tilde F(\tilde V)
$$
We omit proving the compatibility condition; hints: if $V' \to V$, $\hat x'$,
and $x'$ witness the compatibility between $u'$ and $\hat x$, then one
sets $\tilde V' = V' \times_V \tilde V$ which comes with morphsms
$\tilde{\hat x}'$ and $\tilde x'$ and show this works.
The image of $\tilde \xi$ under the transformation $\tilde F \to F$
is the restriction of $\xi$ to $\tilde V$.

\medskip\noindent
Proof of surjectivity: hard case and reduction to affine $W$, part 3.
By our choice of $\tilde W \to W$, there is an affine open
$\tilde V_{open} \subset \tilde V$ (we're running out of notation)
whose image in $V$ contains our chosen point $v \in V$.
Now by the case studied in the next paragraph and the remarks made
earlier, we can descend $\tilde \xi|_{\tilde V_{open}}$
to some element $\tilde \xi_\lambda$ of $\tilde F$ over
$\tilde V_{\lambda, open}$
for some \'etale morphism $\tilde V_{\lambda, open} \to V_\lambda$
whose base change to $V$ is $\tilde V_{open}$.
Applying the transformation of functors $\tilde F \to F$
we obtain the element of $F(\tilde V_{\lambda, open})$
we were looking for. This reduces us to the case discussed
in the next paragraph.

\medskip\noindent
Proof of surjectivity: the case of an affine $W$. We have $v \in Z \subset V$
and $W$ is an affine formal algebraic space. Recall that
$$
\xi = (Z, u', \hat x) \in F(V)
$$
We may still replace $V$ by an \'etale neighbourhood of $v$.
In particular we may and do assume $V$ and $V_\lambda$ are affine.

\medskip\noindent
Proof of surjectivity: descending $Z$. We can find a $\lambda$
and a closed subscheme $Z_\lambda \subset V_\lambda$ such that
$Z$ is the base change of $Z_\lambda$ to $V$. See
Limits, Lemma \ref{limits-lemma-descend-finite-presentation}.
Warning: we don't know (and in general it won't be true)
that $Z_\lambda$ is a reduced closed subscheme of $V_\lambda$.
For $\mu \geq \lambda$ denote $Z_\mu \subset V_\mu$ the scheme theoretic
inverse image in $V_\mu$. We will use below that
$Z = \lim_{\mu \geq \lambda} Z_\mu$ as schemes.

\medskip\noindent
Proof of surjectivity: descending $u'$.
Since $U'$ is locally of finite type over $S$
we may assume after increasing $\lambda$
that there exists a morphism
$u'_\lambda : V_\lambda \setminus Z_\lambda \to U'$
whose restriction to $V \setminus Z$ is $u'$.
See Limits, Proposition
\ref{limits-proposition-characterize-locally-finite-presentation}.
For $\mu \geq \lambda$ we will denote $u'_\mu$ the restriction
of $u'_\lambda$ to $V_\mu \setminus Z_\mu$.

\medskip\noindent
Proof of surjectivity: descending a witness.
Let $V' \to V$, $\hat x'$, and $x'$ witness the compatibility between
$u'$ and $\hat x$. Using the same references as above we may assume
(after increasing $\lambda$) that there exists a morphism
$V'_\lambda \to V_\lambda$ of finite type whose base change to $V$
is $V' \to V$. After increasing $\lambda$ we may assume
$V'_\lambda \to V_\lambda$ is proper
(Limits, Lemma \ref{limits-lemma-eventually-proper}).
Next, we may assume $V'_\lambda \to V_\lambda$ is an isomorphism
over $V_\lambda \setminus Z_\lambda$
(Limits, Lemma \ref{limits-lemma-descend-isomorphism}).
Next, we may assume there is a morphism $x'_\lambda : V'_\lambda \to X'$
whose restriction to $V'$ is $x'$.
Increasing $\lambda$ again we may assume $x'_\lambda$
agrees with $u'_\lambda$ over $V_\lambda \setminus Z_\lambda$.
For $\mu \geq \lambda$ we denote
$V'_\mu$ and $x'_\mu$ the base change of $V'_\lambda$ and the
restriction of $x'_\lambda$.

\medskip\noindent
Proof of surjectivity: algebra.
Write $W = \text{Spf}(B)$, $V = \Spec(A)$, and
for $\mu \geq \lambda$ write $V_\mu = \Spec(A_\mu)$.
Denote $I_\mu \subset A_\mu$ and $I \subset A$
the ideals cutting out $Z_\mu$ and $Z$.
Then $I_\lambda A_\mu = I_\mu$ and $I_\lambda A = I$.
The morphism $\hat x$ determines and is determined by a
continuous ring map
$$
(\hat x)^\sharp : B \longrightarrow A^\wedge
$$
where $A^\wedge$ is the $I$-adic completion of $A$.
To finish the proof we have to show that this map descends to a map into
$A_\mu^\wedge$ for some sufficiently large $\mu$ where $A_\mu^\wedge$
is the $I_\mu$-adic completion of $A_\mu$.
This is a nontrivial fact; Artin writes in his paper
\cite{ArtinII}: ``Since the data (3.5) involve $I$-adic completions,
which do not commute with direct limits, the verification is somewhat
delicate. It is an algebraic analogue of a convergence proof in analysis.''

\medskip\noindent
Proof of surjectivity: algebra, more rings.
Let us denote
$$
C_\mu = \Gamma(V'_\mu, \mathcal{O})
\quad\text{and}\quad
C = \Gamma(V', \mathcal{O})
$$
Observe that $A \to C$ and $A_\mu \to C_\mu$ are finite ring maps as
$V' \to V$ and $V'_\mu \to V_\mu$ are proper morphisms, see
Cohomology of Spaces, Lemma
\ref{spaces-cohomology-lemma-proper-over-affine-cohomology-finite}.
Since $V = \lim V_\mu$ and $V' = \lim V'_\mu$
we have
$$
A = \colim A_\mu
\quad\text{and}\quad
C = \colim C_\mu
$$
by Limits, Lemma \ref{limits-lemma-descend-section}\footnote{We don't
know that $C_\mu = C_\lambda \otimes_{A_\lambda} A_\mu$ as the
various morphisms aren't flat.}. For an element
$a \in I$, resp.\ $a \in I_\mu$ the maps $A_a \to C_a$,
resp.\ $(A_\mu)_a \to (C_\mu)_a$ are isomorphisms by flat base change
(Cohomology of Spaces, Lemma
\ref{spaces-cohomology-lemma-flat-base-change-cohomology}).
Hence the kernel and cokernel of $A \to C$ is supported
on $V(I)$ and similarly for $A_\mu \to C_\mu$.
We conclude the kernel and cokernel of $A \to C$
are annihilated by a power of $I$ and the kernel and cokernel of
$A_\mu \to C_\mu$ are annihilated by a power of $I_\mu$, see
Algebra, Lemma \ref{algebra-lemma-Noetherian-power-ideal-kills-module}.

\medskip\noindent
Proof of surjectivity: algebra, more ring maps.
Denote $Z_n \subset V$ the $n$th infinitesimal
neighbourhood of $Z$ and denote $Z_{\mu, n} \subset V_\mu$
the $n$th infinitesimal neighbourhood of $Z_\mu$.
By the theorem on formal functions
(Cohomology of Spaces, Theorem
\ref{spaces-cohomology-theorem-formal-functions})
we have
$$
C^\wedge = \lim_n H^0(V' \times_V Z_n, \mathcal{O})
\quad\text{and}\quad
C_\mu^\wedge =
\lim_n H^0(V'_\mu \times_{V_\mu} Z_{\mu, n}, \mathcal{O})
$$
where $C^\wedge$ and $C_\mu^\wedge$ are the completion with
respect to $I$ and $I_\mu$.
Combining the completion of the morphism
$x'_\mu : V'_\mu \to X'$ with the morphism $g : X'_{/T'} \to W$ we obtain
$$
g \circ x'_{\mu, /Z_\mu} :
V'_{\mu, /Z_\mu} = \colim V_\mu' \times_{V_\mu} Z_{\mu, n}
\longrightarrow
W
$$
and hence by the description of the completion
$C_\mu^\wedge$ above we obtain a continuous ring homomorphism
$$
(g \circ x'_{\mu, /Z_\mu})^\sharp : B \longrightarrow C_\mu^\wedge
$$
The fact that $V' \to V$, $\hat x'$, $x'$ witnesses the compatibility
between $u'$ and $\hat x$ implies the
commutativity of the following diagram
$$
\xymatrix{
C_\mu^\wedge \ar[r] &
C^\wedge \\
B \ar[u]^{(g \circ x'_{\mu, /Z_\mu})^\sharp} \ar[r]^{(\hat x)^\sharp} &
A^\wedge \ar[u]
}
$$

\medskip\noindent
Proof of surjectivity: more algebra arguments. Recall that the finite
$A$-modules $\Ker(A \to C)$ and $\Coker(A \to C)$ are annihilated
by a power of $I$ and similarly the finite $A_\mu$-modules
$\Ker(A_\mu \to C_\mu)$ and $\Coker(A_\mu \to C_\mu)$ are annihilated
by a power of $I_\mu$. This implies that these modules are
equal to their completions. Since $I$-adic completion on the category of
finite $A$-modules is exact (see
Algebra, Section \ref{algebra-section-completion-noetherian})
it follows that we have
$$
\Coker(A^\wedge \to C^\wedge) = \Coker(A \to C)
$$
and similarly for kernels and for the maps $A_\mu \to C_\mu$.
Of course we also have
$$
\Ker(A \to C) = \colim \Ker(A_\mu \to C_\mu)
\quad\text{and}\quad
\Coker(A \to C) = \colim \Coker(A_\mu \to C_\mu)
$$
Recall that $S = \Spec(R)$ is affine. All of the ring maps
above are $R$-algebra homomorphisms as all of the morphisms
are morphisms over $S$. By
Algebraization of Formal Spaces, Lemma
\ref{restricted-lemma-Noetherian-finite-type-red}
we see that $B$ is topologically of finite type over $R$.
Say $B$ is topologically generated by $b_1, \ldots, b_n$.
Pick some $\mu$ (for example $\lambda$) and consider
the elements
$$
\text{images of }
(g \circ x'_{\mu, /Z_\mu})^\sharp(b_1)
, \ldots,
(g \circ x'_{\mu, /Z_\mu})^\sharp(b_n)
\text{ in }\Coker(A_\mu \to C_\mu)
$$
The image of these elements in $\Coker(\alpha)$ are zero
by the commutativity of the square above. Since
$\Coker(A \to C) = \colim \Coker(A_\mu \to C_\mu)$ and these
cokernels are equal to their completions
we see that after increasing $\mu$ we may assume these
images are all zero. This means that the continuous
homomorphism $(g \circ x'_{\mu, /Z_\mu})^\sharp$ has image contained
in $\Im(A_\mu \to C_\mu)$.
Choose elements $a_{\mu, j} \in (A_\mu)^\wedge$ mapping to
$(g \circ x'_{\mu, /Z_\mu})^\sharp(b_1)$ in $(C_\mu)^\wedge$.
Then $a_{\mu, j} \in A_\mu^\wedge$ and $(\hat x)^\sharp(b_j) \in A^\wedge$
map to the same element of $C^\wedge$ by the commutativity of the
square above. Since
$\Ker(A \to C) = \colim \Ker(A_\mu \to C_\mu)$ and these kernels
are equal to their completions, we may after increasing
$\mu$ adjust our choices of $a_{\mu, j}$ such that
the image of $a_{\mu, j}$ in $A^\wedge$ is equal to $(\hat x)^\sharp(b_j)$.

\medskip\noindent
Proof of surjectivity: final algebra arguments.
Let $\mathfrak b \subset B$ be the ideal of topologically nilpotent
elements. Let $J \subset R[x_1, \ldots, x_n]$ be the ideal
consisting of those $h(x_1, \ldots, x_n)$ such that
$h(b_1, \ldots, b_n) \in \mathfrak b$. Then we get a continuous
surjection of topological $R$-algebras
$$
\Phi : R[x_1, \ldots, x_n]^\wedge \longrightarrow B,\quad
x_j \longmapsto b_j
$$
where the completion on the left hand side is with respect to $J$.
Since $R[x_1, \ldots, x_n]$ is Noetherian we can choose
generators $h_1, \ldots, h_m$ for $J$. By the commutativity
of the square above we see that $h_j(a_{\mu, 1}, \ldots, a_{\mu, n})$ is
an element of $A_\mu^\wedge$ whose image in $A^\wedge$ is contained
in $IA^\wedge$. Namely, the ring map $(\hat x)^\sharp$ is continuous
and $IA^\wedge$ is the ideal of topological nilpotent elements
of $A^\wedge$ because $A^\wedge/IA^\wedge = A/I$ is reduced.
(See Algebra, Section \ref{algebra-section-completion-noetherian}
for results on completion in Noetherian rings.)
Since $A/I = \colim A_\mu/I_\mu$ we conclude that after increasing
$\mu$ we may assume $h_j(a_{\mu, 1}, \ldots, a_{\mu, n})$ is in
$I_\mu A_\mu^\wedge$. In particular the elements
$h_j(a_{\mu, 1}, \ldots, a_{\mu, n})$ of $A_\mu^\wedge$
are topologically nilpotent in $A_\mu^\wedge$.
Thus we obtain a continuous $R$-algebra homomorphism
$$
\Psi : R[x_1, \ldots, x_n]^\wedge \longrightarrow A_\mu^\wedge,\quad
x_j \longmapsto a_{\mu, j}
$$
In order to conclude what we want, we need to see if $\Ker(\Phi)$ is
annihilated by $\Psi$. This may not be true, but we can achieve
this after increasing $\mu$. Indeed, since
$R[x_1, \ldots, x_n]^\wedge$ is Noetherian,
we can choose generators $g_1, \ldots, g_l$ of the ideal
$\Ker(\Phi)$. Then we see that
$$
\Psi(g_1), \ldots, \Psi(g_l) \in
\Ker(A_\mu^\wedge \to C_\mu^\wedge) = \Ker(A_\mu \to C_\mu)
$$
map to zero in $\Ker(A \to C) = \colim \Ker(A_\mu \to C_\mu)$.
Hence increasing $\mu$ as before we get the desired result.

\medskip\noindent
Proof of surjectivity: mopping up. The continuous ring homomorphism
$B \to (A_\mu)^\wedge$ constructed above determines a morphism
$\hat x_\mu : V_{\mu, /Z_\mu} \to W$.
The compatibility of $\hat x_\mu$ and $u'_\mu$ follows
from the fact that the ring map $B \to (A_\mu)^\wedge$
is by construction compatible with the ring map $A_\mu \to C_\mu$.
In fact, the compatibility will be witnessed by the proper morphism
$V'_\mu \to V_\mu$ and the morphisms
$x'_\mu$ and $\hat x'_\mu = x'_{\mu, /Z_\mu}$ we used in
the construction. This finishes the proof.
\end{proof}

\bibliographystyle{plain}
\bibliography{my}
\end{document}
